% Thai Peanut Chicken — Instant Pot
% Save as: InstantPot/Thai_Peanut_Chicken.tex
% Add to main file with: % Thai Peanut Chicken — Instant Pot
% Save as: InstantPot/Thai_Peanut_Chicken.tex
% Add to main file with: % Thai Peanut Chicken — Instant Pot
% Save as: InstantPot/Thai_Peanut_Chicken.tex
% Add to main file with: % Thai Peanut Chicken — Instant Pot
% Save as: InstantPot/Thai_Peanut_Chicken.tex
% Add to main file with: \include{InstantPot/Thai_Peanut_Chicken}
% (Uncomment the \chapter line below only if this is the first recipe in a new chapter.)

%\chapter{Instant Pot}
%\minitoc

\begin{recipebox}{Thai Peanut Chicken}{4}{\textasciitilde30 min (pressure cooker)}

Fast enough for a weeknight, no delay-timer food-safety worries, and the short
cook keeps the breast tender. The rich, thick ingredients go in \emph{after}
pressure cooking to avoid the dreaded ``burn'' warning. Uses 400\,g chicken
bulked out with a tin of chickpeas.

\begin{ingredients}
  \item \textbf{Pressure-cook base}
  \item 400\,g chicken breast, diced into 2--3\,cm cubes
  \item 1 onion, thinly sliced
  \item 3 garlic cloves, crushed
  \item 1 tbsp grated ginger (or 1--2 tsp jar ginger)
  \item \nicefrac{1}{2} cup water or chicken stock
  \item 3 tbsp soy sauce
  \item 1 tbsp fish sauce
  \item 1\nicefrac{1}{2} tsp red curry paste
  \item[] \textbf{Stir in after pressure cooking}
  \item 1 can coconut cream
  \item \nicefrac{1}{4} cup peanut butter
  \item 2 tbsp sweet chilli sauce
  \item 1 tbsp brown sugar or honey
  \item 1 can chickpeas, drained \& rinsed
  \item 2 cups frozen stir-fry vegetables
  \item[] \textbf{To finish}
  \item 2 tbsp peanut butter
  \item 1\nicefrac{1}{2} tbsp lime or lemon juice
  \item 1 tsp toasted sesame oil or rice vinegar (optional)
  \item 1 tbsp cornstarch + 2 tbsp cold water, for slurry (optional)
  \item[] \textbf{To serve}
  \item Rice, crushed peanuts, fresh coriander, extra lime
\end{ingredients}

\begin{method}
  \item \textbf{Sweat the aromatics (optional, worth it).} Set the pot to
        Saut\'e, add a splash of oil, and soften the onion, garlic and ginger
        for 2--3 minutes. Hit Cancel.
  \item \textbf{Build the base.} Add the chicken, water/stock, soy sauce, fish
        sauce and curry paste; stir. This is a thin liquid on purpose --- it
        comes to pressure cleanly and won't burn. Make sure nothing thick is
        stuck to the base.
  \item \textbf{Pressure cook.} Lid on, valve sealed. Pressure Cook / Manual on
        High for \textbf{5 minutes} ($\sim$10 min to come to pressure first).
        When it beeps, \textbf{quick-release} straight away so the chicken stops
        cooking.
  \item \textbf{Stir in sauce \& bulk.} Switch to Saut\'e (low). Stir in the
        coconut cream, \nicefrac{1}{4} cup peanut butter, sweet chilli and sugar
        until smooth, then add the chickpeas and frozen veg. Simmer 4--5 minutes
        until the veg is tender and everything's hot.
  \item \textbf{Finish.} Heat off. Stir in the remaining 2 tbsp peanut butter
        and the lime juice (plus sesame oil/vinegar if using). If thin, stir in
        the cornstarch slurry and bubble a minute on Saut\'e to thicken.
  \item \textbf{Taste \& adjust.} Bland $\rightarrow$ more soy or fish sauce.
        Too rich $\rightarrow$ more lime. Sweeter $\rightarrow$ sweet chilli or
        honey. More heat $\rightarrow$ chilli flakes or sriracha.
  \item \textbf{Serve} over rice, topped with crushed peanuts, coriander and a
        squeeze of lime.
\end{method}

\begin{variationbox}{Slow cooker version}
No burn risk here, so most of it goes in together. Put the sliced onion in the
base and lay the chicken on top --- keep the breasts \emph{whole or in big
chunks} rather than diced so they don't dry out. Pour over everything except the
finishing peanut butter and lime. Slow Cook on Low for 3--4 hours if you're
around (up to 6--8 if you're out all day --- keep the chicken submerged and
shred-and-return it the moment you're home so it drinks the moisture back up).
Stir the chickpeas and frozen veg in for the last 20--30 minutes, then stir the
finishing peanut butter and lime in off the heat. \emph{Avoid the delay-start
timer with raw chicken --- use a shorter cook that drops to Keep Warm instead.}
\end{variationbox}

\begin{notelist}
  \item \textbf{The golden rule:} keep coconut cream, peanut butter, sweet chilli
        and sugar \emph{out} of the pot during pressure cooking --- they scorch
        on the base and trigger the burn error. Thin in first, thick in after.
  \item \textbf{Don't overcook:} 5 minutes high pressure is plenty for 2--3\,cm
        cubes; longer makes breast rubbery. Quick-release to stop it dead.
  \item \textbf{Split the peanut butter} (most in the sauce, a little fresh at
        the end) to keep the flavour bright and nutty rather than flat.
  \item \textbf{Spinach}, if using, goes in right at the end --- it only needs
        the residual heat.
\end{notelist}

\end{recipebox}
% (Uncomment the \chapter line below only if this is the first recipe in a new chapter.)

%\chapter{Instant Pot}
%\minitoc

\begin{recipebox}{Thai Peanut Chicken}{4}{\textasciitilde30 min (pressure cooker)}

Fast enough for a weeknight, no delay-timer food-safety worries, and the short
cook keeps the breast tender. The rich, thick ingredients go in \emph{after}
pressure cooking to avoid the dreaded ``burn'' warning. Uses 400\,g chicken
bulked out with a tin of chickpeas.

\begin{ingredients}
  \item \textbf{Pressure-cook base}
  \item 400\,g chicken breast, diced into 2--3\,cm cubes
  \item 1 onion, thinly sliced
  \item 3 garlic cloves, crushed
  \item 1 tbsp grated ginger (or 1--2 tsp jar ginger)
  \item \nicefrac{1}{2} cup water or chicken stock
  \item 3 tbsp soy sauce
  \item 1 tbsp fish sauce
  \item 1\nicefrac{1}{2} tsp red curry paste
  \item[] \textbf{Stir in after pressure cooking}
  \item 1 can coconut cream
  \item \nicefrac{1}{4} cup peanut butter
  \item 2 tbsp sweet chilli sauce
  \item 1 tbsp brown sugar or honey
  \item 1 can chickpeas, drained \& rinsed
  \item 2 cups frozen stir-fry vegetables
  \item[] \textbf{To finish}
  \item 2 tbsp peanut butter
  \item 1\nicefrac{1}{2} tbsp lime or lemon juice
  \item 1 tsp toasted sesame oil or rice vinegar (optional)
  \item 1 tbsp cornstarch + 2 tbsp cold water, for slurry (optional)
  \item[] \textbf{To serve}
  \item Rice, crushed peanuts, fresh coriander, extra lime
\end{ingredients}

\begin{method}
  \item \textbf{Sweat the aromatics (optional, worth it).} Set the pot to
        Saut\'e, add a splash of oil, and soften the onion, garlic and ginger
        for 2--3 minutes. Hit Cancel.
  \item \textbf{Build the base.} Add the chicken, water/stock, soy sauce, fish
        sauce and curry paste; stir. This is a thin liquid on purpose --- it
        comes to pressure cleanly and won't burn. Make sure nothing thick is
        stuck to the base.
  \item \textbf{Pressure cook.} Lid on, valve sealed. Pressure Cook / Manual on
        High for \textbf{5 minutes} ($\sim$10 min to come to pressure first).
        When it beeps, \textbf{quick-release} straight away so the chicken stops
        cooking.
  \item \textbf{Stir in sauce \& bulk.} Switch to Saut\'e (low). Stir in the
        coconut cream, \nicefrac{1}{4} cup peanut butter, sweet chilli and sugar
        until smooth, then add the chickpeas and frozen veg. Simmer 4--5 minutes
        until the veg is tender and everything's hot.
  \item \textbf{Finish.} Heat off. Stir in the remaining 2 tbsp peanut butter
        and the lime juice (plus sesame oil/vinegar if using). If thin, stir in
        the cornstarch slurry and bubble a minute on Saut\'e to thicken.
  \item \textbf{Taste \& adjust.} Bland $\rightarrow$ more soy or fish sauce.
        Too rich $\rightarrow$ more lime. Sweeter $\rightarrow$ sweet chilli or
        honey. More heat $\rightarrow$ chilli flakes or sriracha.
  \item \textbf{Serve} over rice, topped with crushed peanuts, coriander and a
        squeeze of lime.
\end{method}

\begin{variationbox}{Slow cooker version}
No burn risk here, so most of it goes in together. Put the sliced onion in the
base and lay the chicken on top --- keep the breasts \emph{whole or in big
chunks} rather than diced so they don't dry out. Pour over everything except the
finishing peanut butter and lime. Slow Cook on Low for 3--4 hours if you're
around (up to 6--8 if you're out all day --- keep the chicken submerged and
shred-and-return it the moment you're home so it drinks the moisture back up).
Stir the chickpeas and frozen veg in for the last 20--30 minutes, then stir the
finishing peanut butter and lime in off the heat. \emph{Avoid the delay-start
timer with raw chicken --- use a shorter cook that drops to Keep Warm instead.}
\end{variationbox}

\begin{notelist}
  \item \textbf{The golden rule:} keep coconut cream, peanut butter, sweet chilli
        and sugar \emph{out} of the pot during pressure cooking --- they scorch
        on the base and trigger the burn error. Thin in first, thick in after.
  \item \textbf{Don't overcook:} 5 minutes high pressure is plenty for 2--3\,cm
        cubes; longer makes breast rubbery. Quick-release to stop it dead.
  \item \textbf{Split the peanut butter} (most in the sauce, a little fresh at
        the end) to keep the flavour bright and nutty rather than flat.
  \item \textbf{Spinach}, if using, goes in right at the end --- it only needs
        the residual heat.
\end{notelist}

\end{recipebox}
% (Uncomment the \chapter line below only if this is the first recipe in a new chapter.)

%\chapter{Instant Pot}
%\minitoc

\begin{recipebox}{Thai Peanut Chicken}{4}{\textasciitilde30 min (pressure cooker)}

Fast enough for a weeknight, no delay-timer food-safety worries, and the short
cook keeps the breast tender. The rich, thick ingredients go in \emph{after}
pressure cooking to avoid the dreaded ``burn'' warning. Uses 400\,g chicken
bulked out with a tin of chickpeas.

\begin{ingredients}
  \item \textbf{Pressure-cook base}
  \item 400\,g chicken breast, diced into 2--3\,cm cubes
  \item 1 onion, thinly sliced
  \item 3 garlic cloves, crushed
  \item 1 tbsp grated ginger (or 1--2 tsp jar ginger)
  \item \nicefrac{1}{2} cup water or chicken stock
  \item 3 tbsp soy sauce
  \item 1 tbsp fish sauce
  \item 1\nicefrac{1}{2} tsp red curry paste
  \item[] \textbf{Stir in after pressure cooking}
  \item 1 can coconut cream
  \item \nicefrac{1}{4} cup peanut butter
  \item 2 tbsp sweet chilli sauce
  \item 1 tbsp brown sugar or honey
  \item 1 can chickpeas, drained \& rinsed
  \item 2 cups frozen stir-fry vegetables
  \item[] \textbf{To finish}
  \item 2 tbsp peanut butter
  \item 1\nicefrac{1}{2} tbsp lime or lemon juice
  \item 1 tsp toasted sesame oil or rice vinegar (optional)
  \item 1 tbsp cornstarch + 2 tbsp cold water, for slurry (optional)
  \item[] \textbf{To serve}
  \item Rice, crushed peanuts, fresh coriander, extra lime
\end{ingredients}

\begin{method}
  \item \textbf{Sweat the aromatics (optional, worth it).} Set the pot to
        Saut\'e, add a splash of oil, and soften the onion, garlic and ginger
        for 2--3 minutes. Hit Cancel.
  \item \textbf{Build the base.} Add the chicken, water/stock, soy sauce, fish
        sauce and curry paste; stir. This is a thin liquid on purpose --- it
        comes to pressure cleanly and won't burn. Make sure nothing thick is
        stuck to the base.
  \item \textbf{Pressure cook.} Lid on, valve sealed. Pressure Cook / Manual on
        High for \textbf{5 minutes} ($\sim$10 min to come to pressure first).
        When it beeps, \textbf{quick-release} straight away so the chicken stops
        cooking.
  \item \textbf{Stir in sauce \& bulk.} Switch to Saut\'e (low). Stir in the
        coconut cream, \nicefrac{1}{4} cup peanut butter, sweet chilli and sugar
        until smooth, then add the chickpeas and frozen veg. Simmer 4--5 minutes
        until the veg is tender and everything's hot.
  \item \textbf{Finish.} Heat off. Stir in the remaining 2 tbsp peanut butter
        and the lime juice (plus sesame oil/vinegar if using). If thin, stir in
        the cornstarch slurry and bubble a minute on Saut\'e to thicken.
  \item \textbf{Taste \& adjust.} Bland $\rightarrow$ more soy or fish sauce.
        Too rich $\rightarrow$ more lime. Sweeter $\rightarrow$ sweet chilli or
        honey. More heat $\rightarrow$ chilli flakes or sriracha.
  \item \textbf{Serve} over rice, topped with crushed peanuts, coriander and a
        squeeze of lime.
\end{method}

\begin{variationbox}{Slow cooker version}
No burn risk here, so most of it goes in together. Put the sliced onion in the
base and lay the chicken on top --- keep the breasts \emph{whole or in big
chunks} rather than diced so they don't dry out. Pour over everything except the
finishing peanut butter and lime. Slow Cook on Low for 3--4 hours if you're
around (up to 6--8 if you're out all day --- keep the chicken submerged and
shred-and-return it the moment you're home so it drinks the moisture back up).
Stir the chickpeas and frozen veg in for the last 20--30 minutes, then stir the
finishing peanut butter and lime in off the heat. \emph{Avoid the delay-start
timer with raw chicken --- use a shorter cook that drops to Keep Warm instead.}
\end{variationbox}

\begin{notelist}
  \item \textbf{The golden rule:} keep coconut cream, peanut butter, sweet chilli
        and sugar \emph{out} of the pot during pressure cooking --- they scorch
        on the base and trigger the burn error. Thin in first, thick in after.
  \item \textbf{Don't overcook:} 5 minutes high pressure is plenty for 2--3\,cm
        cubes; longer makes breast rubbery. Quick-release to stop it dead.
  \item \textbf{Split the peanut butter} (most in the sauce, a little fresh at
        the end) to keep the flavour bright and nutty rather than flat.
  \item \textbf{Spinach}, if using, goes in right at the end --- it only needs
        the residual heat.
\end{notelist}

\end{recipebox}
% (Uncomment the \chapter line below only if this is the first recipe in a new chapter.)

%\chapter{Instant Pot}
%\minitoc

\begin{recipebox}{Thai Peanut Chicken}{4}{\textasciitilde30 min (pressure cooker)}

Fast enough for a weeknight, no delay-timer food-safety worries, and the short
cook keeps the breast tender. The rich, thick ingredients go in \emph{after}
pressure cooking to avoid the dreaded ``burn'' warning. Uses 400\,g chicken
bulked out with a tin of chickpeas.

\begin{ingredients}
  \item \textbf{Pressure-cook base}
  \item 400\,g chicken breast, diced into 2--3\,cm cubes
  \item 1 onion, thinly sliced
  \item 3 garlic cloves, crushed
  \item 1 tbsp grated ginger (or 1--2 tsp jar ginger)
  \item \nicefrac{1}{2} cup water or chicken stock
  \item 3 tbsp soy sauce
  \item 1 tbsp fish sauce
  \item 1\nicefrac{1}{2} tsp red curry paste
  \item[] \textbf{Stir in after pressure cooking}
  \item 1 can coconut cream
  \item \nicefrac{1}{4} cup peanut butter
  \item 2 tbsp sweet chilli sauce
  \item 1 tbsp brown sugar or honey
  \item 1 can chickpeas, drained \& rinsed
  \item 2 cups frozen stir-fry vegetables
  \item[] \textbf{To finish}
  \item 2 tbsp peanut butter
  \item 1\nicefrac{1}{2} tbsp lime or lemon juice
  \item 1 tsp toasted sesame oil or rice vinegar (optional)
  \item 1 tbsp cornstarch + 2 tbsp cold water, for slurry (optional)
  \item[] \textbf{To serve}
  \item Rice, crushed peanuts, fresh coriander, extra lime
\end{ingredients}

\begin{method}
  \item \textbf{Sweat the aromatics (optional, worth it).} Set the pot to
        Saut\'e, add a splash of oil, and soften the onion, garlic and ginger
        for 2--3 minutes. Hit Cancel.
  \item \textbf{Build the base.} Add the chicken, water/stock, soy sauce, fish
        sauce and curry paste; stir. This is a thin liquid on purpose --- it
        comes to pressure cleanly and won't burn. Make sure nothing thick is
        stuck to the base.
  \item \textbf{Pressure cook.} Lid on, valve sealed. Pressure Cook / Manual on
        High for \textbf{5 minutes} ($\sim$10 min to come to pressure first).
        When it beeps, \textbf{quick-release} straight away so the chicken stops
        cooking.
  \item \textbf{Stir in sauce \& bulk.} Switch to Saut\'e (low). Stir in the
        coconut cream, \nicefrac{1}{4} cup peanut butter, sweet chilli and sugar
        until smooth, then add the chickpeas and frozen veg. Simmer 4--5 minutes
        until the veg is tender and everything's hot.
  \item \textbf{Finish.} Heat off. Stir in the remaining 2 tbsp peanut butter
        and the lime juice (plus sesame oil/vinegar if using). If thin, stir in
        the cornstarch slurry and bubble a minute on Saut\'e to thicken.
  \item \textbf{Taste \& adjust.} Bland $\rightarrow$ more soy or fish sauce.
        Too rich $\rightarrow$ more lime. Sweeter $\rightarrow$ sweet chilli or
        honey. More heat $\rightarrow$ chilli flakes or sriracha.
  \item \textbf{Serve} over rice, topped with crushed peanuts, coriander and a
        squeeze of lime.
\end{method}

\begin{variationbox}{Slow cooker version}
No burn risk here, so most of it goes in together. Put the sliced onion in the
base and lay the chicken on top --- keep the breasts \emph{whole or in big
chunks} rather than diced so they don't dry out. Pour over everything except the
finishing peanut butter and lime. Slow Cook on Low for 3--4 hours if you're
around (up to 6--8 if you're out all day --- keep the chicken submerged and
shred-and-return it the moment you're home so it drinks the moisture back up).
Stir the chickpeas and frozen veg in for the last 20--30 minutes, then stir the
finishing peanut butter and lime in off the heat. \emph{Avoid the delay-start
timer with raw chicken --- use a shorter cook that drops to Keep Warm instead.}
\end{variationbox}

\begin{notelist}
  \item \textbf{The golden rule:} keep coconut cream, peanut butter, sweet chilli
        and sugar \emph{out} of the pot during pressure cooking --- they scorch
        on the base and trigger the burn error. Thin in first, thick in after.
  \item \textbf{Don't overcook:} 5 minutes high pressure is plenty for 2--3\,cm
        cubes; longer makes breast rubbery. Quick-release to stop it dead.
  \item \textbf{Split the peanut butter} (most in the sauce, a little fresh at
        the end) to keep the flavour bright and nutty rather than flat.
  \item \textbf{Spinach}, if using, goes in right at the end --- it only needs
        the residual heat.
\end{notelist}

\end{recipebox}